% =====================================================================
% SAH-RAG: Architecture and Implementation Report
% A Technical Deep-Dive into the Structure-Aware Hierarchical Hybrid
% Retrieval System (rag_system)
% =====================================================================
\documentclass[11pt,a4paper]{article}

\usepackage[utf8]{inputenc}
\usepackage[T1]{fontenc}
\usepackage{lmodern}
\usepackage[margin=1in]{geometry}
\usepackage{amsmath,amssymb,amsfonts}
\usepackage{booktabs}
\usepackage{multirow}
\usepackage{longtable}
\usepackage{array}
\usepackage{tabularx}
\usepackage{siunitx}
\usepackage{graphicx}
\usepackage{xcolor}
\usepackage{hyperref}
\usepackage{url}
\usepackage{enumitem}
\usepackage{listings}
\usepackage{fancyhdr}
\usepackage{titlesec}
\usepackage{caption}
\usepackage{subcaption}
\usepackage{microtype}
\usepackage{parskip}
\usepackage{algorithm}
\usepackage{algpseudocode}

\hypersetup{
  colorlinks=true,
  linkcolor=blue!50!black,
  citecolor=blue!50!black,
  urlcolor=blue!50!black,
  pdftitle={SAH-RAG Architecture and Implementation Report},
  pdfauthor={RAG Architecture Readiness Program},
}

\sisetup{round-mode=places,round-precision=4,detect-weight=true}

\definecolor{codebg}{RGB}{248,248,248}
\definecolor{codekw}{RGB}{0,90,160}
\definecolor{codestr}{RGB}{160,50,20}
\definecolor{codecom}{RGB}{110,110,110}

\lstdefinelanguage{python}{
  language=Python,
  morekeywords={self,True,False,None},
  keywordstyle=\color{codekw}\bfseries,
  commentstyle=\color{codecom}\itshape,
  stringstyle=\color{codestr},
}
\lstset{
  basicstyle=\ttfamily\footnotesize,
  backgroundcolor=\color{codebg},
  frame=single,
  rulecolor=\color{black!20},
  breaklines=true,
  showstringspaces=false,
  captionpos=b,
  columns=fullflexible,
  keepspaces=true,
  numbers=left,
  numberstyle=\tiny\color{black!40},
}

\pagestyle{fancy}
\fancyhf{}
\rhead{\small My\_RAG --- Architecture \& Implementation}
\lhead{\small rag\_system Technical Reference}
\cfoot{\thepage}

\title{
  \vspace{-1.5cm}
  \rule{\linewidth}{0.4pt}\\[0.3cm]
  {\Large\bfseries My\_RAG (\texttt{rag\_system}):}\\[0.15cm]
  {\large Architecture, Design Rationale, and Implementation\\ of a Structure-Aware Hierarchical Hybrid Retriever}\\[0.2cm]
  \rule{\linewidth}{0.4pt}\\[0.4cm]
  \normalsize RAG Architecture Readiness Program\\
  \small Technical Implementation Report
}
\author{}
\date{\today}

\begin{document}

\maketitle
\thispagestyle{empty}

% =====================================================================
\begin{abstract}
\noindent
We present a complete technical description of \emph{My\_RAG}
(\texttt{rag\_system}), the project-owned reference retrieval
architecture of the RAG Architecture Readiness Program. The
architecture is a \emph{structure-aware hierarchical hybrid} retriever
that treats a canonical document as a composition of deterministic
semantic parents and child chunks, indexes child chunks in a Qdrant
collection that carries both a dense BGE-M3 vector and a learned-sparse
BGE-M3 vector, retrieves candidates independently from both
representations, and fuses their rankings with Reciprocal Rank Fusion
before aggregating back to the parent and document levels. The present
document describes the architecture at implementation granularity: the
exact data models and their field-level constraints, the algorithms
that produce each artifact, the deterministic identity schemes that
make every layer content-addressable, and the invariants that each
layer is required to maintain. We do not report experimental results;
those are documented separately in the program's comparison and
operating-point reports. The report is organized as a code-adjacent
walkthrough of the nine subpackages of \texttt{rag\_system}
(\texttt{canonical}, \texttt{parsing}, \texttt{chunking},
\texttt{embedding}, \texttt{storage}, \texttt{retrieval},
\texttt{evaluation}, \texttt{artifacts}, and the top-level
orchestration modules), with cross-references to the specific
classes, functions, and configuration fields that define each stage.
The goal of this document is to make every design decision and every
implementation detail of My\_RAG legible to a reader who has not
seen the source code.
\end{abstract}

\vfill
\noindent\textbf{Keywords.} retrieval-augmented generation;
hierarchical chunking; semantic parents; learned-sparse retrieval;
reciprocal rank fusion; deterministic artifacts; content-addressable
identity; Pydantic data models; Qdrant hybrid collections.

\newpage
\tableofcontents
\newpage

% =====================================================================
\section{Introduction}\label{sec:intro}

\subsection{Purpose of this document}
This document is a \emph{technical implementation report} for My\_RAG,
the project-owned retrieval architecture of the RAG Architecture
Readiness Program. It is the companion document to:

\begin{itemize}[leftmargin=1.5em,itemsep=0.15em]
  \item the My\_RAG baseline report, which establishes the
    architecture's place in the program and reports its retrieval
    quality against the shared evaluation protocol;
  \item the My\_RAG optimal-operation study, which characterizes
    the architecture's parameter-response behavior and freezes a
    validated operating configuration;
  \item the LongRAG vs.\ My\_RAG comparison report, which places
    My\_RAG's numbers beside an external baseline.
\end{itemize}

What those reports do \emph{not} do is describe \emph{how} the
architecture is built. The present report is that description. It
contains no experimental results, no metric tables, and no
cross-architecture comparisons. Its purpose is to make every design
decision and every implementation detail of My\_RAG legible to a
reader who has not seen the source code.

\subsection{Scope and audience}
The intended reader is a retrieval engineer or researcher who wants
to understand:

\begin{itemize}[leftmargin=1.5em,itemsep=0.15em]
  \item what data models My\_RAG uses and what invariants those
    models enforce;
  \item what algorithms transform the input at each stage and what
    those algorithms guarantee;
  \item why each design decision was made and what alternatives
    were considered;
  \item how reproducibility is enforced at every layer.
\end{itemize}

The report is organized as a code-adjacent walkthrough. Each section
corresponds to a subpackage or module of the codebase, presents the
data models that the module exposes, describes the algorithms that
operate on those models, and states the invariants the module is
responsible for maintaining.

\subsection{How to read this report}
Section~\ref{sec:principles} states the five design principles that
govern the architecture. Section~\ref{sec:architecture} gives a
system-level overview. Sections~\ref{sec:canonical}
through~\ref{sec:evaluation} descend into each pipeline layer in
order. Section~\ref{sec:artifacts} covers the artifact infrastructure
that every layer depends on. Section~\ref{sec:determinism}
summarizes the determinism and reproducibility guarantees that the
architecture as a whole provides. Section~\ref{sec:packages}
catalogs the package structure and entry points.

% =====================================================================
\section{Design Principles}\label{sec:principles}

My\_RAG is governed by five design principles. Every implementation
decision in this report traces back to one or more of them.

\subsection{Preserve source structure}
The architecture must never silently discard or blur the structural
organization of a source document. A document's sections, clauses,
and paragraphs are not an accident of formatting; they are evidence
of how the document's author intended the material to be read. Any
retrieval unit that crosses a section boundary mixes evidence that
the author deliberately separated.

This principle is implemented at two levels. First, the structural
parser produces \emph{semantic parents} that are guaranteed to be
contiguous, non-overlapping spans of the normalized document text.
Second, the child chunker never emits a chunk that crosses a parent
boundary. Together, these two invariants mean that every child chunk
in the system is a fragment of exactly one semantic parent, and every
semantic parent is a fragment of exactly one canonical document.

\subsection{Retrieve fine, rank coarse}
The granularity at which evidence is best matched to a query is not
the granularity at which the user wants results. Retrieval works
well on small, semantically focused units: a query is compared to a
fragment of text, and the fragment's meaning is measured by its
embedding. But the user wants to see documents, not fragments.

My\_RAG therefore retrieves at the child-chunk level and ranks at
the document level. Between the two, an intermediate parent level
aggregates child evidence into coherent parent scores. The two-step
aggregation (child $\to$ parent $\to$ document) preserves the fine
granularity of the match while returning results at the granularity
of interest.

\subsection{Combine complementary signals}
Dense and learned-sparse retrieval have different failure modes.
Dense retrieval captures semantic similarity but can miss exact
lexical matches in technical or domain-specific queries. Learned-sparse
retrieval captures lexical evidence but cannot generalize across
paraphrases. Neither alone is best on every query.

The architecture retrieves independently from both representations
and fuses their rankings. Because both representations are produced
by the same model (BGE-M3), they share the same vocabulary and
tokenizer, which makes the fusion cleaner than fusing two independently
trained signals.

The fusion uses Reciprocal Rank Fusion (RRF), which is rank-based
rather than score-based. This is a deliberate choice: RRF does not
require the two representations' scores to be on the same scale, and
it does not privilege whichever representation happens to produce
larger raw scores.

\subsection{Be deterministic everywhere}
Every identity, every tie-break, every serialization format, and
every aggregation rule is a deterministic function of the input. Two
executions of the same pipeline on the same inputs produce byte-identical
artifacts (modulo floating-point nondeterminism in the embedder, which
is isolated to the embedding layer).

Determinism matters for three reasons. First, it makes the architecture
auditable: any artifact can be regenerated from its inputs and
verified against a checksum. Second, it makes the architecture
resumable: an interrupted run can be restarted without producing
inconsistent state. Third, it makes the architecture fair to compare
against other architectures: no cross-run variance can be attributed
to the architecture's own behavior.

\subsection{Enforce invariants loudly}
Where a property is required for correctness, the architecture
enforces it at the earliest possible layer and raises an exception
if it is violated. Silent data loss is treated as a bug. Examples
include:

\begin{itemize}[leftmargin=1.5em,itemsep=0.15em]
  \item a \texttt{SemanticParent} whose text does not equal the
    slice of the canonical normalized text at its character span is
    rejected;
  \item a \texttt{ChildChunk} whose normalized span is not contained
    inside its parent's normalized span is rejected;
  \item a chunked document whose chunk indices are not contiguous
    from zero is rejected;
  \item a Qdrant ingestion batch that contains a duplicate point ID
    is rejected;
  \item an embedding shard whose records do not cover the full set
    of source chunks is rejected.
\end{itemize}

The principle is that any invariant violation is a signal that
something upstream has gone wrong, and it is better to fail at the
violating layer than to propagate corrupted data downstream.

% =====================================================================
\section{System Architecture}\label{sec:architecture}

\subsection{Dataflow}
Figure~\ref{fig:pipeline} shows the architecture's end-to-end
dataflow. The offline portion reads canonical documents and produces
a Qdrant hybrid collection. The online portion reads a query and
produces a ranked document list.

\begin{figure}[h]
\centering
\small
\begin{tabular}{l}
\hline
\textbf{Offline}\\
\hline
\texttt{CanonicalDocument}\\
$\downarrow$ \texttt{DeterministicStructuralParser.parse()} \\
\texttt{ParsedDocument} containing a tuple of \texttt{SemanticParent} \\
$\downarrow$ \texttt{HierarchicalChildChunker.chunk\_document()} \\
tuple of \texttt{ChildChunk} \\
$\downarrow$ BGE-M3 (dense + learned sparse) \\
\texttt{ChunkEmbeddingRecord} per child chunk \\
$\downarrow$ \texttt{ingest\_embedding\_records()} \\
Qdrant hybrid collection (named vectors \texttt{dense} and \texttt{sparse}) \\
\hline
\textbf{Online}\\
\hline
query text\\
$\downarrow$ BGE-M3 (dense + learned sparse) \\
query dense vector, query sparse vector \\
$\downarrow$ \texttt{retrieve\_dense()} / \texttt{retrieve\_sparse()} / \texttt{retrieve\_hybrid\_rrf()} \\
ranked \texttt{RetrievedChunk} \\
$\downarrow$ \texttt{aggregate\_chunks\_to\_parents()} \\
ranked \texttt{RetrievedParent} \\
$\downarrow$ \texttt{resolve\_parents\_to\_documents()} \\
ranked \texttt{RetrievedDocument} \\
\hline
\end{tabular}
\caption{My\_RAG end-to-end dataflow.}
\label{fig:pipeline}
\end{figure}

\subsection{Layer responsibilities}
Table~\ref{tab:layers} summarizes the responsibility of each layer and
the module that implements it.

\begin{table}[h]
\centering
\caption{Subpackage responsibilities.}
\label{tab:layers}
\small
\begin{tabular}{p{3.2cm} p{11cm}}
\toprule
\textbf{Subpackage} & \textbf{Responsibility} \\
\midrule
\texttt{canonical} & Typed canonical document and example models; loaders; path resolution. \\
\texttt{parsing} & Deterministic structural parser producing semantic parents. \\
\texttt{chunking} & Parent-aware token chunker producing child chunks. \\
\texttt{embedding} & Dense and learned-sparse embedding record models; shard I/O; catalog. \\
\texttt{storage} & Qdrant collection creation and hybrid ingestion. \\
\texttt{retrieval} & Dense, sparse, and RRF retrieval; parent aggregation; document resolution. \\
\texttt{evaluation} & Ranked-retrieval metrics; evaluation case construction; manifest writing. \\
\texttt{artifacts} & Base Pydantic model; deterministic identities; sidecar manifests; atomic I/O. \\
\texttt{generation} & (present but optional; disabled in retrieval-only execution). \\
\bottomrule
\end{tabular}
\end{table}

% =====================================================================
\section{Canonical Input Layer}\label{sec:canonical}

The \texttt{canonical} subpackage defines the input data models and
the streaming loaders that read them from disk. It contains five
modules: \texttt{models.py}, \texttt{loaders.py}, \texttt{paths.py},
\texttt{jsonl.py}, and \texttt{inspection.py}.

\subsection{Data models}
The two principal models are \texttt{CanonicalDocument} and
\texttt{CanonicalExample}.

\subsubsection{CanonicalDocument}
\begin{lstlisting}[language=python]
class CanonicalDocument(CanonicalBaseModel):
    doc_id: str
    dataset: str
    text_raw: str
    text_normalized: str
    title: str | None = None
    source_path: str | None = None
    source_row: int | str | None = None
    metadata: JsonObject = Field(default_factory=dict)
\end{lstlisting}

The model has \texttt{extra="allow"} (inherited from
\texttt{CanonicalBaseModel}), which means that additional fields
present in a canonical file are preserved rather than rejected. It
has \texttt{frozen=True}, so instances are immutable once created.
Field validators require that \texttt{doc_id} and \texttt{dataset}
are non-empty after stripping whitespace.

Both \texttt{text\_raw} and \texttt{text\_normalized} are carried.
The raw text is the exact source bytes; the normalized text is the
result of whitespace normalization applied during canonical dataset
construction. The parser and chunker operate exclusively on the
normalized text, but the raw text is preserved so that character
offsets in the normalized text can be mapped back to the raw text
in principle (though the mapping is only exact when the two texts
are identical; the parser handles the case where they are not by
setting \texttt{raw\_char\_span} to \texttt{None} for the affected
artifacts).

\subsubsection{CanonicalExample, CanonicalQuery, CanonicalGroundTruth}
An example bundles a query and its ground truth:
\begin{lstlisting}[language=python]
class CanonicalQuery(CanonicalBaseModel):
    query_id: str
    dataset: str
    question: str
    task_type: str | None = None
    split: str | None = None
    doc_ids: list[str] = Field(default_factory=list)
    metadata: JsonObject = Field(default_factory=dict)

class CanonicalGroundTruth(CanonicalBaseModel):
    ground_truth_type: GroundTruthType
    dataset: str
    query_id: str
    answers: list[Any] = Field(default_factory=list)
    relevant_ids: list[str] = Field(default_factory=list)
    choice: Any | None = None
    spans: list[Any] = Field(default_factory=list)
    metadata: JsonObject = Field(default_factory=dict)

class CanonicalExample(CanonicalBaseModel):
    query: CanonicalQuery
    ground_truth: CanonicalGroundTruth
    supplied_context: Any | None = None
    options: Any | None = None
    metadata: JsonObject = Field(default_factory=dict)
\end{lstlisting}

The \texttt{GroundTruthType} enum has five values:
\textsc{Composite}, \textsc{SpanSet}, \textsc{FreeText},
\textsc{MultipleChoice}, and \textsc{Deterministic}. The enum is used
to signal to downstream consumers which ground-truth fields are
populated; the program's retrieval-only evaluator uses
\texttt{relevant\_ids} regardless of type.

\texttt{CanonicalExample} enforces an identity consistency check at
\texttt{model\_post\_init}: the query's \texttt{query\_id} must equal
the ground truth's \texttt{query\_id}, and the query's
\texttt{dataset} must equal the ground truth's \texttt{dataset}. This
prevents constructing an example whose query and ground truth refer
to different benchmarks.

\subsection{Loaders}
The \texttt{loaders.py} module provides two public streaming
functions:

\begin{lstlisting}[language=python]
def iter_canonical_documents(
    dataset: str | CanonicalDatasetPaths,
) -> Iterator[CanonicalDocument]

def iter_canonical_examples(
    dataset: str | CanonicalDatasetPaths,
) -> Iterator[CanonicalExample]
\end{lstlisting}

Both functions are true iterators: they stream records from JSONL,
validating each record against its Pydantic model before yielding it.
A record that fails validation raises a
\texttt{CanonicalModelValidationError} that carries the file path,
the record number within the file, the model name, and the underlying
Pydantic \texttt{ValidationError}. This structured error is what
allows the pipeline to abort loudly and specifically when a canonical
file is corrupted.

The loaders do not load the full dataset into memory. They are safe
to use on the largest canonical files in the suite, which exceed
\SI{800}{MB}.

\subsection{Paths}
The \texttt{paths.py} module resolves a dataset name to a
\texttt{CanonicalDatasetPaths} dataclass containing:

\begin{itemize}[leftmargin=1.5em,itemsep=0.15em]
  \item \texttt{documents} --- the path to \texttt{documents.jsonl};
  \item \texttt{examples} --- the path to \texttt{examples.jsonl};
  \item \texttt{validation} --- the path to \texttt{validation.json};
  \item \texttt{manifest} --- the path to \texttt{manifest.json}.
\end{itemize}

The root directory for a dataset is determined by the
\texttt{get\_canonical\_root()} function, which reads from a
project-level configuration or environment variable. This indirection
is what allows the canonical suite to be relocated without code
changes.

% =====================================================================
\section{Structural Parsing}\label{sec:parsing}

The \texttt{parsing} subpackage converts a \texttt{CanonicalDocument}
into a \texttt{ParsedDocument}, which is a structured representation
containing zero or more \texttt{SemanticParent} objects. The parser
is deterministic: the same input text always produces the same parent
sequence, with the same parent identities.

\subsection{The \texttt{SemanticParent} model}
\begin{lstlisting}[language=python]
class SemanticParent(ArtifactBaseModel):
    parent_id: str
    doc_id: str
    dataset: str

    parent_parent_id: str | None = None
    hierarchy_level: int = Field(ge=0)
    sibling_index: int = Field(ge=0)

    title: str | None = None
    text_normalized: str
    text_raw: str | None = None

    normalized_char_span: OffsetSpan
    raw_char_span: OffsetSpan | None = None

    token_span: OffsetSpan | None = None
    token_count: int | None = Field(default=None, ge=0)

    start_boundary: tuple[BoundaryEvidence, ...] = ()
    end_boundary: tuple[BoundaryEvidence, ...] = ()

    parser_id: str
    parser_version: str
    parser_config_hash: str

    metadata: JsonObject = Field(default_factory=dict)
\end{lstlisting}

The model carries three kinds of information:

\begin{enumerate}[leftmargin=1.5em,itemsep=0.15em]
  \item \textbf{Identity}: \texttt{parent\_id}, \texttt{doc\_id},
    \texttt{dataset}. All three must be non-empty.
  \item \textbf{Content}: \texttt{text\_normalized} and
    \texttt{text\_raw}, together with the character spans
    (\texttt{normalized\_char\_span}, \texttt{raw\_char\_span}) that
    locate the parent inside its canonical document.
  \item \textbf{Provenance}: \texttt{start\_boundary} and
    \texttt{end\_boundary} are tuples of \texttt{BoundaryEvidence}
    objects, each carrying a \texttt{method} (from the
    \texttt{BoundaryMethod} enum), an optional \texttt{confidence},
    and optional \texttt{details}. This design allows a boundary to
    be justified by multiple independent pieces of evidence.
\end{enumerate}

The model validators enforce the following invariants:

\begin{itemize}[leftmargin=1.5em,itemsep=0.15em]
  \item a parent cannot be its own parent
    (\texttt{parent\_parent\_id != parent\_id});
  \item if the normalized character span is non-empty, the normalized
    text must also be non-empty.
\end{itemize}

The \texttt{BoundaryMethod} enum enumerates the ways a boundary can
be justified. The parser's own methods are
\textsc{DocumentStart}, \textsc{DocumentEnd}, \textsc{ExplicitHeading},
and \textsc{DocumentPreamble}; other methods exist in the enum for
compatibility with parsers that use different signals.

\subsection{The parsing algorithm}
The parser's entry point is:

\begin{lstlisting}[language=python]
class DeterministicStructuralParser:
    def __init__(
        self,
        config: DeterministicParserConfig | None = None,
    ) -> None: ...

    def parse(
        self,
        document: CanonicalDocument,
    ) -> ParsedDocument: ...
\end{lstlisting}

The \texttt{parse()} method executes four steps: segmentation,
heading detection, grouping, and parent construction. Each step is
described below.

\subsubsection{Step 1: Paragraph segmentation}
The function \texttt{iter\_paragraph\_spans()} in
\texttt{parsing/segmentation.py} splits the normalized text into
blocks using the regular expression
\texttt{(?:\textbackslash r?\textbackslash n[ \textbackslash t]*)\{2,\}},
which matches two or more newlines (with optional intervening
whitespace). Each block is trimmed of leading and trailing
whitespace, and its \texttt{start} and \texttt{end} offsets in the
original text are preserved exactly.

Empty blocks are discarded. The result is a tuple of \texttt{TextBlock}
dataclasses, each carrying a \texttt{block\_index}, \texttt{text},
\texttt{start}, and \texttt{end}. The \texttt{TextBlock} dataclass
validates that \texttt{len(text) == end - start} and that the text is
non-empty.

\subsubsection{Step 2: Heading detection}
The function \texttt{detect\_heading\_reason()} in
\texttt{parsing/headings.py} classifies a block as a heading using a
conservative three-tier scheme. The function returns a
\texttt{HeadingReason} enum value (\textsc{Numbered}, \textsc{Named},
or \textsc{Uppercase}) or \texttt{None}.

Before any of the tier-specific rules are applied, the block must
pass a set of basic filters:
\begin{itemize}[leftmargin=1.5em,itemsep=0.15em]
  \item the text must be non-empty;
  \item it must not contain a newline;
  \item it must be at most \texttt{maximum\_heading\_characters} long
    (default \num{120});
  \item it must be at most \texttt{maximum\_heading\_words} words
    (default \num{14});
  \item it must contain at least
    \texttt{minimum\_heading\_alphabetic\_characters} alphabetic
    characters (default \num{4});
  \item it must not contain URL-like or equation-like patterns;
  \item it must not consist solely of numeric/table fragments;
  \item it must contain at most
    \texttt{maximum\_heading\_terminal\_punctuation} terminal
    punctuation marks (default \num{1}).
\end{itemize}

After the basic filters, the tier-specific rules are applied in
order:

\begin{itemize}[leftmargin=1.5em,itemsep=0.15em]
  \item \textbf{Named headings}: the block matches the
    \texttt{\_NAMED\_HEADING} regex, which accepts English words
    such as \texttt{chapter}, \texttt{section}, \texttt{article},
    \texttt{part}, \texttt{title}, \texttt{appendix}, \texttt{annex},
    \texttt{exhibit}, \texttt{schedule}, as well as Greek equivalents.
    The match is case-insensitive.
  \item \textbf{Numbered headings}: the block starts with a numbered
    prefix matching \texttt{\_SECTION\_NUMBER} (patterns like
    \texttt{1.}, \texttt{1.2.3}, \texttt{(1)}, \texttt{IV.}, or
    \texttt{A.}) and additionally passes
    \texttt{\_looks\_like\_numbered\_heading()}, which rejects
    single-token blocks, blocks whose remainder has no alphabetic
    content, and blocks ending with sentence-like punctuation.
  \item \textbf{Uppercase headings}: the block's alphabetic
    characters are at least \num{90}\% uppercase and it does not
    consist of a single very short word.
\end{itemize}

The three-tier approach is deliberately conservative. The cost of
misclassifying a heading is that a section boundary is placed where
the author did not intend one; the cost of missing a heading is that
two intended sections are merged into one parent. Both are less
harmful than over-eager heading detection, which is why the rules
require multiple affirmative signals.

\subsubsection{Step 3: Grouping}
The grouping step partitions the block sequence into sections, each
beginning at a heading block (or at the first block of the document
if the document starts with a preamble). Concretely:

\begin{enumerate}[leftmargin=1.5em,itemsep=0.15em]
  \item Identify all blocks classified as headings.
  \item If there are no headings, return the entire block sequence
    as a single group.
  \item Otherwise, if the first heading is not the first block,
    emit the preamble as its own group.
  \item For each heading index $i$, emit the group consisting of
    blocks $i$ through the next heading index minus one (or through
    the end of the sequence for the final heading).
\end{enumerate}

If \texttt{merge\_small\_structural\_groups} is enabled (the default),
the grouping step then merges any group whose character length is
below \texttt{minimum\_parent\_characters} (default \num{20}) into
its successor, or into its predecessor if it has no successor. This
prevents the parser from emitting tiny orphan parents that are too
small to be useful retrieval units.

\subsubsection{Step 4: Parent construction}
For each group, the parser builds one \texttt{SemanticParent}. The
parent's character span is from the first block's \texttt{start} to
the last block's \texttt{end}. The parent's normalized text is the
slice of the canonical normalized text at that span. If the canonical
document's \texttt{text\_raw} equals its
\texttt{text\_normalized} (which is the common case), the raw span
and raw text are also populated.

The parent's \texttt{parent\_id} is built by
\texttt{build\_parent\_id()}, which takes the dataset, document ID,
parser configuration hash, and parent index and produces a
content-addressable identifier of the form
\texttt{parent-\{dataset\}-\{index:06d\}-\{fingerprint\}}, where the
fingerprint is the first twelve hex characters of the SHA-256 hash of
a canonical JSON serialization of the identity tuple.

The parent's boundary evidence is set as follows:
\begin{itemize}[leftmargin=1.5em,itemsep=0.15em]
  \item the first parent's start boundary is
    \textsc{DocumentStart};
  \item a parent whose first block is a heading has a start boundary
    of \textsc{ExplicitHeading};
  \item any other parent has a start boundary of
    \textsc{DocumentPreamble};
  \item the last parent's end boundary is \textsc{DocumentEnd};
  \item any other parent's end boundary is \textsc{ParagraphBreak}.
\end{itemize}

\subsection{Invariants maintained by the parser}
The parser's output is a \texttt{ParsedDocument} whose
\texttt{parents} tuple satisfies:

\begin{enumerate}[leftmargin=1.5em,itemsep=0.15em]
  \item each parent's text equals the exact slice of the canonical
    normalized text at the parent's character span;
  \item parent spans are disjoint and their union covers the
    document's non-whitespace content;
  \item parent IDs are unique within the document;
  \item each parent's \texttt{doc\_id} and \texttt{dataset} match the
    document's;
  \item each parent's parser identity fields match the parser that
    produced it.
\end{enumerate}

The parser's \texttt{parse()} method also computes a
\texttt{source\_fingerprint} that combines the document ID, dataset,
raw text, and normalized text into a single SHA-256 hash. This
fingerprint is used by downstream layers to detect any drift between
the document that was parsed and the document that is currently on
disk.

% =====================================================================
\section{Hierarchical Child Chunking}\label{sec:chunking}

The \texttt{chunking} subpackage converts a \texttt{ParsedDocument}
into a tuple of \texttt{ChildChunk} objects. Like the parser, it is
deterministic: the same parsed document with the same chunker
configuration always produces the same chunks with the same
identities.

\subsection{The \texttt{ChildChunk} model}
\begin{lstlisting}[language=python]
class ChildChunk(ArtifactBaseModel):
    artifact_type: str = "child_chunk"
    artifact_version: str = "1.0"

    chunk_id: str
    parent_id: str
    doc_id: str
    dataset: str

    chunk_index: int = Field(ge=0)
    parent_chunk_index: int = Field(ge=0)

    text_normalized: str
    text_raw: str | None = None

    normalized_char_span: OffsetSpan
    parent_normalized_char_span: OffsetSpan

    raw_char_span: OffsetSpan | None = None
    token_span: OffsetSpan | None = None
    token_count: int = Field(ge=0)

    parser_id: str
    parser_version: str
    parser_config_hash: str

    chunker_id: str
    chunker_version: str
    chunker_config_hash: str

    tokenizer_id: str
    chunk_size: int = Field(gt=0)
    chunk_overlap: int = Field(ge=0)

    metadata: JsonObject = Field(default_factory=dict)
\end{lstlisting}

The model distinguishes between two indices:

\begin{itemize}[leftmargin=1.5em,itemsep=0.15em]
  \item \texttt{chunk\_index} is the chunk's position in the
    document-wide sequence of chunks, starting at zero and
    incrementing contiguously across all parents;
  \item \texttt{parent\_chunk\_index} is the chunk's position within
    its parent, starting at zero for each new parent.
\end{itemize}

Two character spans are carried: \texttt{normalized\_char\_span} is
the chunk's absolute span in the canonical normalized text, and
\texttt{parent\_normalized\_char\_span} is the containing parent's
span. Both are needed: the former for exact source reconstruction,
the latter for verifying that the chunk is contained inside its
parent.

The model validators enforce:

\begin{itemize}[leftmargin=1.5em,itemsep=0.15em]
  \item all identity fields are non-empty;
  \item \texttt{chunk\_overlap < chunk\_size};
  \item the chunk's normalized span is contained inside the parent's
    normalized span;
  \item a chunk with a positive token count has non-empty text;
  \item if \texttt{token\_span} is set, its length equals
    \texttt{token\_count}.
\end{itemize}

\subsection{Tokenizer abstraction}
The chunker does not depend on any specific tokenizer. It depends on
the \texttt{OffsetTokenizer} protocol defined in
\texttt{chunking/tokenization.py}:

\begin{lstlisting}[language=python]
@runtime_checkable
class OffsetTokenizer(Protocol):
    @property
    def tokenizer_id(self) -> str: ...

    @property
    def special_tokens_count(self) -> int: ...

    def tokenize_with_offsets(
        self,
        text: str,
    ) -> tuple[TokenSpan, ...]: ...
\end{lstlisting}

The protocol requires that the tokenizer produce \texttt{TokenSpan}
objects with exact character offsets. This is the property that
allows the chunker to produce chunks whose text is exactly the slice
of the source text covered by the chunk's token window.

Two concrete implementations are provided:

\begin{itemize}[leftmargin=1.5em,itemsep=0.15em]
  \item \texttt{RegexOffsetTokenizer} is a dependency-free
    tokenizer used for tests. It splits on the pattern
    \texttt{\textbackslash w+|[^\textbackslash w\textbackslash s]}
    with Unicode support.
  \item \texttt{HuggingFaceOffsetTokenizer} adapts a fast Hugging
    Face tokenizer to the protocol. It requires that the wrapped
    tokenizer expose \texttt{is\_fast = True} and produce
    \texttt{offset\_mapping}; it raises \texttt{ValueError} if either
    precondition fails.
\end{itemize}

The Hugging Face adapter's \texttt{tokenizer\_id} is a structured
string of the form
\texttt{huggingface:\{model\_id\}@\{revision\}:fast}, which pins the
tokenizer identity to both the model and the revision used. This is
what makes the chunker's \texttt{config\_hash} stable across
different runs that happen to use the same tokenizer.

The adapter's \texttt{tokenize\_with\_offsets()} method encodes the
text with \texttt{add\_special\_tokens=False}, so the resulting token
sequence contains only content tokens; the special tokens that BGE-M3
would add when embedding are accounted for separately via
\texttt{special\_tokens\_count}. Any token whose offset has
\texttt{end <= start} (which can happen for special tokens or
zero-width markers) is skipped.

\subsection{Chunker configuration}
\begin{lstlisting}[language=python]
class ChildChunkerConfig(BaseModel):
    model_config = ConfigDict(extra="forbid", frozen=True)

    chunker_id: str = "hierarchical_token_chunker"
    chunker_version: str = "0.3.0"

    chunk_size: int = Field(default=256, gt=0)
    chunk_overlap: int = Field(default=32, ge=0)

    minimum_final_chunk_tokens: int | None = Field(
        default=None, gt=0,
    )

    preserve_parent_boundaries: bool = True
    include_empty_parents: bool = False
\end{lstlisting}

The configuration has two derived properties. The
\texttt{content\_chunk\_size()} method computes the number of content
tokens available per chunk, which is \texttt{chunk\_size} minus the
tokenizer's \texttt{special\_tokens\_count}. The
\texttt{effective\_minimum\_final\_chunk\_tokens()} method returns
\texttt{minimum\_final\_chunk\_tokens} if it is set, or half the
content chunk size (floored to an integer) otherwise.

The model validator enforces \texttt{chunk\_overlap < chunk\_size}
and \texttt{minimum\_final\_chunk\_tokens <= chunk\_size}, and it
rejects any attempt to set
\texttt{preserve\_parent\_boundaries = False}. The last constraint is
what guarantees that the chunker will never produce a chunk that
crosses a parent boundary.

\subsection{The chunking algorithm}
The chunker's entry point is:

\begin{lstlisting}[language=python]
class HierarchicalChildChunker:
    def __init__(
        self,
        *,
        tokenizer: OffsetTokenizer,
        config: ChildChunkerConfig | None = None,
    ) -> None: ...

    def chunk_document(
        self,
        parsed_document: ParsedDocument,
    ) -> tuple[ChildChunk, ...]: ...
\end{lstlisting}

The constructor computes the chunker's \texttt{config\_hash} from the
configuration fields plus the tokenizer's identity and special-token
count. This hash is what makes the chunker's output content-addressable
in the same sense as the parser's.

The \texttt{chunk\_document()} method iterates over the parsed
document's parents and calls \texttt{\_chunk\_parent()} on each. Each
parent's chunks are appended to a document-wide list, and the
\texttt{global\_start\_index} is advanced by the number of chunks
produced for that parent. After all parents have been processed, the
method calls \texttt{\_validate\_document\_chunks()} to verify the
document-level invariants (unique chunk IDs, contiguous global
indices, valid parent references, matching document and dataset
identifiers).

\subsubsection{Chunking one parent}
The \texttt{\_chunk\_parent()} method:

\begin{enumerate}[leftmargin=1.5em,itemsep=0.15em]
  \item tokenizes the parent's normalized text using the tokenizer;
  \item if the tokenizer returns no tokens, either returns a single
    empty chunk (if \texttt{include\_empty\_parents} is set) or
    returns an empty tuple;
  \item otherwise, calls \texttt{build\_token\_windows()} to determine
    the token windows for the parent's content;
  \item for each window, calls \texttt{\_build\_chunk()} to construct
    a \texttt{ChildChunk}.
\end{enumerate}

\subsubsection{Token windows}
The function \texttt{build\_token\_windows()} in
\texttt{chunking/windows.py} is where the chunking policy is
implemented. It takes the total token count, the content chunk size,
the chunk overlap, and the minimum final chunk size, and returns a
tuple of \texttt{TokenWindow} half-open intervals.

The algorithm has three characteristics that matter:

\begin{enumerate}[leftmargin=1.5em,itemsep=0.15em]
  \item \textbf{It never produces a window shorter than the minimum
    final chunk size as a trailing residue.} When the natural
    next window would produce a final chunk shorter than the
    minimum, the algorithm instead slides the window backward so
    that it is anchored at the end of the parent. This
    \emph{anchored final window} guarantees that the last chunk is
    large enough to be useful.
  \item \textbf{It never produces two identical adjacent windows.}
    The final deduplication step removes any window that is
    identical to its predecessor, which can happen in edge cases
    where the anchor logic produces a window that coincides with
    the previous window.
  \item \textbf{It is deterministic.} Given the same inputs, it
    always produces the same sequence of windows. This is what makes
    the chunker's output reproducible.
\end{enumerate}

\subsubsection{Building a chunk}
The \texttt{\_build\_chunk()} method takes the window's first and
last \texttt{TokenSpan} and uses their character offsets to compute
the chunk's text and span:

\begin{enumerate}[leftmargin=1.5em,itemsep=0.15em]
  \item \texttt{relative\_start} and \texttt{relative\_end} are the
    first and last token's offsets relative to the parent's
    normalized text;
  \item \texttt{text\_normalized} is the slice of the parent's
    normalized text at those offsets;
  \item \texttt{normalized\_char\_span} is computed by adding the
    parent's \texttt{normalized\_char\_span.start} to the relative
    offsets;
  \item if the parent has a raw span and raw text, the chunk's raw
    span and raw text are computed analogously.
\end{enumerate}

The chunk's \texttt{chunk\_id} is built by
\texttt{build\_chunk\_id()}, which combines the dataset, document ID,
parent ID, chunker configuration hash, and \texttt{parent\_chunk\_index}.
The identifier has the form
\texttt{chunk-\{dataset\}-\{index:06d\}-\{fingerprint\}}.

The chunk's metadata block records a set of computed values that
downstream layers use for verification and reporting:
\begin{itemize}[leftmargin=1.5em,itemsep=0.15em]
  \item the relative character start and end offsets within the
    parent;
  \item whether this is the first and/or last chunk of its parent;
  \item the model-input token budget;
  \item the content-token capacity per chunk;
  \item the special-token count;
  \item the estimated model-input token count for the chunk
    (\texttt{token\_count + special\_tokens\_count}).
\end{itemize}

\subsection{Invariants maintained by the chunker}
The chunker's output is a tuple of \texttt{ChildChunk} objects that
satisfy:

\begin{enumerate}[leftmargin=1.5em,itemsep=0.15em]
  \item chunk IDs are unique within the document;
  \item global chunk indices are contiguous from zero;
  \item every chunk's parent ID is a known parent of the parsed
    document;
  \item every chunk's document ID and dataset match the parsed
    document's;
  \item every chunk's normalized span is contained inside its
    parent's;
  \item every chunk's text equals the exact slice of its parent's
    text at the chunk's relative span.
\end{enumerate}

% =====================================================================
\section{Embedding Layer}\label{sec:embedding}

The \texttt{embedding} subpackage defines the models that carry
embeddings, the shard-writer that produces embedding artifacts, and
the catalog that validates a complete embedding run.

\subsection{Embedding models}
Three models are defined in \texttt{embedding/models.py}:

\subsubsection{EmbeddingKind and VectorPrecision}
Two enums classify embeddings:

\begin{lstlisting}[language=python]
class EmbeddingKind(StrEnum):
    DENSE = "dense"
    LEARNED_SPARSE = "learned_sparse"
    HYBRID = "hybrid"

class VectorPrecision(StrEnum):
    FLOAT16 = "float16"
    FLOAT32 = "float32"
\end{lstlisting}

\texttt{EmbeddingKind} distinguishes dense vectors, learned-sparse
vectors, and shards that carry both. \texttt{VectorPrecision}
distinguishes FP16 and FP32 dense vectors.

\subsubsection{SparseVector}
\begin{lstlisting}[language=python]
class SparseVector(ArtifactBaseModel):
    indices: tuple[int, ...]
    values: tuple[float, ...]
\end{lstlisting}

The validators require that \texttt{indices} and \texttt{values} have
the same length, that all indices are non-negative, that indices are
unique, that indices are sorted in ascending order, and that all
values are finite. The sorted-indices requirement is what allows the
storage layer to serialize sparse vectors as a compact pair of arrays
without additional sorting at write time.

\subsubsection{ChunkEmbeddingRecord}
\begin{lstlisting}[language=python]
class ChunkEmbeddingRecord(ArtifactBaseModel):
    artifact_type: str = "chunk_embedding"
    artifact_version: str = "1.0"

    chunk_id: str
    parent_id: str
    doc_id: str
    dataset: str

    model_id: str
    model_revision: str
    tokenizer_id: str
    embedding_config_hash: str

    dense_vector: tuple[float, ...] | None = None
    sparse_vector: SparseVector | None = None

    dense_dimension: int | None = Field(default=None, gt=0)
    dense_precision: VectorPrecision | None = None
    dense_normalized: bool | None = None

    source_chunk_sha256: str | None = None
\end{lstlisting}

The record requires at least one of the two representations to be
present. If the dense vector is present, all three of
\texttt{dense\_dimension}, \texttt{dense\_precision}, and
\texttt{dense\_normalized} are required, and the vector length must
equal \texttt{dense\_dimension}. If the dense vector is absent, all
three of those fields must also be absent. If
\texttt{source\_chunk\_sha256} is set, it must be exactly sixty-four
hexadecimal characters.

\subsubsection{EmbeddingShardHeader}
\begin{lstlisting}[language=python]
class EmbeddingShardHeader(ArtifactBaseModel):
    artifact_type: str = "embedding_shard_header"
    artifact_version: str = "1.0"

    shard_id: str
    shard_index: int = Field(ge=0)

    dataset: str
    model_id: str
    model_revision: str
    tokenizer_id: str
    embedding_config_hash: str

    embedding_kind: EmbeddingKind

    dense_dimension: int | None = Field(default=None, gt=0)
    dense_precision: VectorPrecision | None = None
    dense_normalized: bool | None = None

    maximum_records: int = Field(gt=0)
\end{lstlisting}

The header carries the identity and configuration of one shard. Its
validators require that dense metadata be present if and only if the
embedding kind is \textsc{Dense} or \textsc{Hybrid}.

\subsection{Embedding configuration}
\begin{lstlisting}[language=python]
class EmbeddingConfig(BaseModel):
    model_config = ConfigDict(extra="forbid", frozen=True)

    embedder_id: str = "bge_m3_embedder"
    embedder_version: str = "0.1.0"

    model_id: str = "BAAI/bge-m3"
    model_revision: str

    tokenizer_id: str

    embedding_kind: EmbeddingKind = EmbeddingKind.HYBRID

    dense_dimension: int | None = Field(default=1024, gt=0)
    dense_precision: VectorPrecision | None = VectorPrecision.FLOAT32
    dense_normalized: bool | None = True

    device: str = "cpu"
    use_fp16: bool = False

    batch_size: int = Field(default=2, gt=0)
    maximum_model_input_tokens: int = Field(default=8192, gt=0)
    shard_maximum_records: int = Field(default=10000, gt=0)
\end{lstlisting}

The configuration's validators enforce that all identity fields are
non-empty, that dense metadata is present if and only if the kind is
dense or hybrid, and that \texttt{use\_fp16} is not combined with
\texttt{device = "cpu"}. The last check exists because the initial
CPU profile of the pipeline does not enable FP16; a future profile
that does would need to loosen this constraint.

\subsection{Shard I/O}
The \texttt{embedding/shards.py} module provides the
\texttt{write\_embedding\_shard()} function, which is the canonical
way to write embedding records to disk:

\begin{lstlisting}[language=python]
def write_embedding_shard(
    *,
    root: Path,
    header: EmbeddingShardHeader,
    records: Iterable[ChunkEmbeddingRecord],
    source_chunks_path: Path,
    source_chunks_manifest_path: Path | None = None,
    overwrite: bool = False,
    execution_profile: str | None = None,
    experiment_id: str | None = None,
    metadata: dict[str, Any] | None = None,
) -> EmbeddingShardWriteSummary: ...
\end{lstlisting}

The function performs the following operations:

\begin{enumerate}[leftmargin=1.5em,itemsep=0.15em]
  \item Resolves the shard's conventional file paths
    (\texttt{header.json}, \texttt{records.jsonl},
    \texttt{shard.manifest.json}).
  \item Loads the source chunk IDs from the source chunk JSONL file
    to use for membership validation.
  \item Writes the records through an \texttt{AtomicJsonlWriter},
    validating each record against the header and against the source
    chunk ID set.
  \item Writes the header atomically.
  \item Builds and writes a sidecar manifest that records the SHA-256
    of the header, the records file, and every input.
\end{enumerate}

If any step fails, the previously written files are deleted before
the exception propagates. This is what makes the writer atomic: a
failed write leaves no partial artifacts.

The record-writer's validation checks every record against the
header's identity fields (\texttt{dataset}, \texttt{model\_id},
\texttt{model\_revision}, \texttt{tokenizer\_id},
\texttt{embedding\_config\_hash}), against the header's
representation kind (dense vs.\ sparse vs.\ hybrid), and against the
header's dense metadata (dimension, precision, normalization). It also
verifies that the record's \texttt{chunk\_id} is present in the source
chunk ID set.

\subsection{Embedding catalog}
The \texttt{embedding/catalog.py} module defines the
\texttt{EmbeddingRunCatalog} model, which describes a complete
embedding run across multiple shards, and the
\texttt{validate\_embedding\_catalog\_files()} function, which
verifies the catalog against the files on disk.

The catalog's validators enforce:
\begin{itemize}[leftmargin=1.5em,itemsep=0.15em]
  \item shard indices are contiguous from zero;
  \item the sum of shard record counts equals the catalog's total
    record count;
  \item the shards' source ranges are contiguous and cover the full
    record count.
\end{itemize}

The file-level validator additionally verifies that each shard's
header and manifest match the catalog entry's checksums, that each
shard's record count matches the catalog entry's, and that no chunk
ID appears in more than one shard.

% =====================================================================
\section{Storage Layer}\label{sec:storage}

The \texttt{storage} subpackage provides Qdrant collection creation
and hybrid-vector ingestion.

\subsection{Collection creation}
\begin{lstlisting}[language=python]
def create_hybrid_collection(
    *,
    client: QdrantClient,
    collection_name: str,
    dense_dimension: int,
    recreate: bool = False,
) -> None:
\end{lstlisting}

The function creates a Qdrant collection with two named vectors:
\texttt{dense} (the dense BGE-M3 vector, COSINE distance,
unquantized) and \texttt{sparse} (the learned-sparse vector,
in-memory index, unquantized). Both vector names are module-level
constants (\texttt{DENSE\_VECTOR\_NAME = "dense"} and
\texttt{SPARSE\_VECTOR\_NAME = "sparse"}) that are also exported
from the storage package's \texttt{\_\_init\_\_}.

The function refuses to overwrite an existing collection unless
\texttt{recreate=True} is passed explicitly. This is a deliberate
safety measure: accidental collection replacement is a common source
of silent data loss in retrieval systems.

\subsection{Deterministic point identity}
\begin{lstlisting}[language=python]
def build_qdrant_point_id(chunk_id: str) -> str:
    return str(uuid.uuid5(uuid.NAMESPACE_URL, f"rag-chunk:{chunk_id}"))
\end{lstlisting}

Qdrant point IDs are computed from the chunk ID using UUID version 5
with the URL namespace. This makes the point identity a deterministic
function of the chunk identity: the same chunk always maps to the
same point ID, regardless of ingestion order or Qdrant instance. This
is what makes Qdrant ingestion idempotent.

\subsection{Payload construction}
The \texttt{build\_payload()} function constructs the searchable
payload for a Qdrant point. The payload includes:

\begin{itemize}[leftmargin=1.5em,itemsep=0.15em]
  \item the chunk's identity (\texttt{chunk\_id}, \texttt{parent\_id},
    \texttt{doc\_id}, \texttt{dataset});
  \item the chunk's indices (\texttt{chunk\_index},
    \texttt{parent\_chunk\_index});
  \item the chunk's normalized text;
  \item the chunk's character and token spans;
  \item the model and embedding identity (\texttt{model\_id},
    \texttt{model\_revision}, \texttt{embedding\_config\_hash});
  \item the shard ID.
\end{itemize}

The payload is checked before construction: the chunk's and record's
chunk IDs must match, and the record's dense vector length must match
the header's declared dimension.

\subsection{Ingestion}
\begin{lstlisting}[language=python]
def ingest_embedding_records(
    *,
    client: QdrantClient,
    collection_name: str,
    header: EmbeddingShardHeader,
    records: Iterable[ChunkEmbeddingRecord],
    chunks: Iterable[ChildChunk],
    batch_size: int = 64,
) -> QdrantIngestionSummary:
\end{lstlisting}

The function indexes the chunks by chunk ID, iterates over the
records, builds a Qdrant point for each, and upserts points in
batches of \texttt{batch\_size}. It validates that every record
references a known chunk, that no two points in the same batch have
the same ID, and that the header's dense dimension is present (since
hybrid ingestion requires it).

The return value is a \texttt{QdrantIngestionSummary} that records
the collection name, the number of points ingested, the dense
dimension, and both named-vector names.

\subsection{Collection manifest}
The \texttt{QdrantCollectionManifest} model describes a completed
collection. Its validators enforce that all identity fields are
non-empty, that the catalog and configuration hashes are valid
SHA-256 strings, and that the manifest's \texttt{complete} flag is
consistent with its point counts.

% =====================================================================
\section{Retrieval Layer}\label{sec:retrieval}

The \texttt{retrieval} subpackage implements the online retrieval
path: dense retrieval, sparse retrieval, RRF fusion, parent
aggregation, and document resolution.

\subsection{Retrieval models}
Two models are defined in \texttt{retrieval/models.py}:

\subsubsection{RetrievedChunk}
\begin{lstlisting}[language=python]
class RetrievedChunk(ArtifactBaseModel):
    chunk_id: str
    parent_id: str
    doc_id: str
    dataset: str

    rank: int = Field(ge=1)
    score: float

    text_normalized: str
    retrieval_method: str

    payload: dict[str, object] = Field(default_factory=dict)
\end{lstlisting}

The model carries one child chunk with its score and rank. Its
validators require that all identity fields are non-empty.

\subsubsection{RetrievedParent}
\begin{lstlisting}[language=python]
class RetrievedParent(ArtifactBaseModel):
    parent_id: str
    doc_id: str
    dataset: str

    rank: int = Field(ge=1)
    score: float

    best_chunk_rank: int = Field(ge=1)
    matched_chunk_count: int = Field(ge=1)

    chunk_ids: tuple[str, ...]
    child_scores: tuple[float, ...]
\end{lstlisting}

The parent model carries the aggregated score, the best contributing
child's rank, and the count of contributing children, plus the
contributing chunk IDs and their individual scores. Its validators
enforce that the chunk IDs and child scores are aligned, that
\texttt{matched\_chunk\_count} equals the number of contributing
children, and that no chunk ID appears twice.

\subsection{Dense retrieval}
\begin{lstlisting}[language=python]
def retrieve_dense(
    *,
    client: QdrantClient,
    collection_name: str,
    dense_vector: Sequence[float],
    limit: int,
) -> tuple[RetrievedChunk, ...]:
\end{lstlisting}

The function issues a Qdrant query against the collection's
\texttt{dense} named vector, with \texttt{limit} set to the number
of candidates requested. The Qdrant response is converted to a tuple
of \texttt{RetrievedChunk} objects, with the rank assigned
sequentially starting at one.

The conversion function \texttt{\_convert\_points()} validates that
each Qdrant payload has the required fields
(\texttt{chunk\_id}, \texttt{parent\_id}, \texttt{doc\_id},
\texttt{dataset}, \texttt{text\_normalized}) before constructing the
typed result. A payload that is missing any of these fields raises a
\texttt{ValueError} with the names of the missing fields.

\subsection{Sparse retrieval}
\begin{lstlisting}[language=python]
def retrieve_sparse(
    *,
    client: QdrantClient,
    collection_name: str,
    sparse_indices: Sequence[int],
    sparse_values: Sequence[float],
    limit: int,
) -> tuple[RetrievedChunk, ...]:
\end{lstlisting}

The function issues a Qdrant query against the collection's
\texttt{sparse} named vector, passing the query's sparse vector
constructed from the given indices and values. Its return value is
analogous to the dense case.

\subsection{Hybrid RRF retrieval}
\begin{lstlisting}[language=python]
def retrieve_hybrid_rrf(
    *,
    client: QdrantClient,
    collection_name: str,
    dense_vector: Sequence[float],
    sparse_indices: Sequence[int],
    sparse_values: Sequence[float],
    candidate_limit: int,
    result_limit: int,
) -> tuple[RetrievedChunk, ...]:
\end{lstlisting}

The function issues a Qdrant query with two \texttt{Prefetch}
clauses, one against the \texttt{dense} vector and one against the
\texttt{sparse} vector, each with \texttt{limit = candidate\_limit}.
The outer query uses \texttt{FusionQuery(fusion=Fusion.RRF)} and
returns \texttt{result\_limit} results.

The function validates that \texttt{candidate\_limit >= result\_limit},
that \texttt{result\_limit > 0}, and that the sparse indices and
values are aligned in length.

The use of Qdrant's native RRF is deliberate: it keeps the fusion
step inside the database, avoids transferring candidate sets back to
Python for fusion, and guarantees that the fusion is performed on the
full ranked lists rather than on truncated lists.

\subsection{Parent aggregation}
\begin{lstlisting}[language=python]
def aggregate_chunks_to_parents(
    chunks: tuple[RetrievedChunk, ...],
    *,
    additional_child_weight: float = 0.05,
    limit: int | None = None,
) -> tuple[RetrievedParent, ...]:
\end{lstlisting}

The aggregation rule is:

\[
\mathrm{score}(P) = \max_{c \in P} \mathrm{score}(c)
  + w \sum_{c \in P \setminus \{c^\star\}} \mathrm{score}(c)
\]

where $c^\star$ is the highest-scoring child of parent $P$ and
$w = \texttt{additional\_child\_weight}$.

The algorithm:

\begin{enumerate}[leftmargin=1.5em,itemsep=0.15em]
  \item Groups the retrieved chunks by \texttt{parent\_id}.
  \item For each parent, sorts the group's members by rank and by
    descending score, producing two views.
  \item Computes the parent's score as the best score plus the
    weighted sum of the remaining scores.
  \item Records the best chunk's rank as the parent's
    \texttt{best\_chunk\_rank}, and the number of contributing
    children as \texttt{matched\_chunk\_count}.
  \item Sorts parents by descending score, then by best chunk rank,
    then by parent ID.
  \item If \texttt{limit} is set, truncates to the first \texttt{limit}
    parents.
\end{enumerate}

The weight $w$ is a hyperparameter. A value of zero reduces the rule
to best-only; a value of one would sum all children's scores. The
default of \num{0.05} was chosen to credit additional evidence
without over-rewarding parents that simply have more children.

\subsection{Document resolution}
\begin{lstlisting}[language=python]
def resolve_parents_to_documents(
    parents: Sequence[RetrievedParent],
    *,
    query_id: str,
    limit: int | None = None,
) -> tuple[RetrievedDocument, ...]:
\end{lstlisting}

The document resolution rule is: a document's score is the score of
its highest-scoring parent.

The algorithm:

\begin{enumerate}[leftmargin=1.5em,itemsep=0.15em]
  \item Groups parents by \texttt{doc\_id}.
  \item For each document, sorts its parents by rank and by descending
    score, producing two views.
  \item The document's score is the highest parent's score; the
    document's \texttt{best\_parent\_rank} is the best-ranked parent's
    rank.
  \item Accumulates all contributing parent IDs and chunk IDs,
    deduplicated, in the order in which they first appear.
  \item Sorts documents by descending score, then by best parent
    rank, then by document ID.
  \item If \texttt{limit} is set, truncates to the first \texttt{limit}
    documents.
\end{enumerate}

The resulting \texttt{RetrievedDocument} carries the document ID, its
rank, its score, its best parent rank, its contributing parent and
chunk IDs, and a metadata block. The metadata block is validated
against a set of prohibited keys that includes \texttt{generated\_answer},
\texttt{reference\_answer}, \texttt{answer\_text}, \texttt{relevance},
\texttt{judgment}, and similar terms. The purpose of this prohibition
is to make it impossible to accidentally leak ground truth into the
retrieval artifact.

% =====================================================================
\section{Evaluation Layer}\label{sec:evaluation}

The \texttt{evaluation} subpackage defines the ranked-retrieval metric
set, the evaluation case model, the pipeline that writes evaluation
artifacts, and the paired-bootstrap helper.

\subsection{Evaluation case}
\begin{lstlisting}[language=python]
class RetrievalJudgment(ArtifactBaseModel):
    item_id: str
    relevance: int = Field(ge=0)

class RetrievalEvaluationCase(ArtifactBaseModel):
    query_id: str
    query_text: str
    retrieval_method: str
    evaluation_level: str

    judgments: tuple[RetrievalJudgment, ...]
    retrieved_ids: tuple[str, ...]

    metadata: dict[str, object] = Field(default_factory=dict)
\end{lstlisting}

A case bundles a query, its relevance judgments, and the ranked
retrieved IDs for one retrieval method. The validators require that
judgment IDs are unique and that retrieved IDs are unique.

The separation of judgments from the retrieval call is deliberate:
it is what allows the program to enforce that ground truth is
loaded only during evaluation and never during retrieval.

\subsection{Metrics}
The \texttt{evaluation/metrics.py} module implements the standard
ranked-retrieval metric set:

\begin{lstlisting}[language=python]
def recall_at_k(
    *, retrieved_ids: Sequence[str],
    relevance_by_id: Mapping[str, int],
    cutoff: int,
) -> float

def precision_at_k(
    *, retrieved_ids: Sequence[str],
    relevance_by_id: Mapping[str, int],
    cutoff: int,
) -> float

def reciprocal_rank(
    *, retrieved_ids: Sequence[str],
    relevance_by_id: Mapping[str, int],
    cutoff: int,
) -> float

def ndcg_at_k(
    *, retrieved_ids: Sequence[str],
    relevance_by_id: Mapping[str, int],
    cutoff: int,
) -> float
\end{lstlisting}

The metrics are binary recall, fixed-cutoff precision, first-relevant
reciprocal rank, and graded nDCG with gain formula
$2^{\mathrm{rel}} - 1$ and discount $1 / \log_2(\mathrm{rank} + 1)$.

\texttt{evaluate\_case()} evaluates one case at one cutoff,
constructing a \texttt{RetrievalMetrics} record with the query ID,
retrieval method, evaluation level, cutoff, relevant count, retrieved
count, relevant-retrieved count, and the five metric values (recall,
precision, reciprocal rank, nDCG, hit).

\texttt{aggregate\_metrics()} computes macro averages over a sequence
of metrics records that share the same retrieval method, evaluation
level, and cutoff. It rejects mixed batches.

\subsection{Pipeline}
The \texttt{write\_retrieval\_evaluation()} function writes the
artifacts of an evaluation run:

\begin{itemize}[leftmargin=1.5em,itemsep=0.15em]
  \item \texttt{cases.jsonl}: the evaluated cases;
  \item \texttt{metrics.jsonl}: one metrics record per case per
    cutoff;
  \item \texttt{aggregates.jsonl}: one aggregate per (method, level,
    cutoff) group;
  \item \texttt{run.json}: the run identity and metadata;
  \item \texttt{evaluation.manifest.json}: a sidecar manifest with
    checksums and inputs.
\end{itemize}

The run's \texttt{run\_id} is built deterministically from the
dataset ID, evaluation level, retrieval methods, cutoffs, and a
configuration hash.

\subsection{Paired bootstrap}
The \texttt{evaluation/statistics.py} module provides a
paired-bootstrap helper used for uncertainty estimates in the
operating-point study. It is not part of the retrieval pipeline
itself.

% =====================================================================
\section{Artifact Infrastructure}\label{sec:artifacts}

The \texttt{artifacts} subpackage provides the base Pydantic model,
the deterministic identity functions, the sidecar manifest model,
the atomic JSON writer, and the file-fingerprinting helpers that
every other subpackage depends on.

\subsection{Base model}
\begin{lstlisting}[language=python]
class ArtifactBaseModel(BaseModel):
    model_config = ConfigDict(
        extra="forbid",
        frozen=True,
        str_strip_whitespace=False,
    )
\end{lstlisting}

All derived artifacts inherit from \texttt{ArtifactBaseModel}. The
three configurations are:

\begin{itemize}[leftmargin=1.5em,itemsep=0.15em]
  \item \texttt{extra="forbid"} rejects any field not declared on the
    model. This is what catches schema drift early.
  \item \texttt{frozen=True} makes instances immutable.
  \item \texttt{str\_strip\_whitespace=False} preserves the exact
    string content, which matters for character-offset integrity.
\end{itemize}

\subsection{Deterministic identities}
The \texttt{artifacts/identity.py} module provides:

\begin{itemize}[leftmargin=1.5em,itemsep=0.15em]
  \item \texttt{canonical\_json\_bytes()} serializes any
    JSON-compatible value with sorted keys and no whitespace,
    producing a canonical byte representation;
  \item \texttt{sha256\_fingerprint()} hashes the canonical bytes;
  \item \texttt{short\_fingerprint()} returns the first $n$
    hex characters of the hash;
  \item \texttt{normalize\_identifier()} lowercases and
    alphanumeric-normalizes a string for inclusion in an ID;
  \item \texttt{build\_parent\_id()}, \texttt{build\_chunk\_id()},
    \texttt{build\_experiment\_id()}, and
    \texttt{combine\_fingerprints()} construct the specific ID forms
    used by the pipeline layers.
\end{itemize}

The ID construction pattern is consistent across all three:
\begin{enumerate}[leftmargin=1.5em,itemsep=0.15em]
  \item build an identity dictionary with all the inputs;
  \item compute the identity's SHA-256 fingerprint;
  \item concatenate the type prefix, the normalized dataset, the
    human-readable index, and the fingerprint suffix.
\end{enumerate}

\subsection{Sidecar manifest}
\begin{lstlisting}[language=python]
class ArtifactSidecarManifest(ArtifactBaseModel):
    manifest_type: str = "artifact_sidecar_manifest"
    manifest_version: str = "1.0"

    artifact_type: str
    artifact_path: str
    artifact_sha256: str
    artifact_size_bytes: int = Field(ge=0)
    record_count: int = Field(ge=0)

    dataset_id: str | None = None
    experiment_id: str | None = None

    component: ArtifactComponentReference
    inputs: tuple[ArtifactInputReference, ...] = ()

    code_revision: str | None = None
    execution_profile: str | None = None

    created_at: datetime = Field(
        default_factory=lambda: datetime.now(UTC)
    )
    completed: bool = True

    metadata: JsonObject = Field(default_factory=dict)
\end{lstlisting}

The manifest records the artifact's identity, size, SHA-256,
record count, and the component that produced it. It also records
the code revision (from \texttt{get\_git\_revision()}) and an
execution profile string.

\texttt{ArtifactComponentReference} carries the component type,
component ID, version, configuration hash, and the configuration
dictionary itself.

\texttt{ArtifactInputReference} carries the type, path, SHA-256,
size, and optional record count of a consumed input artifact. Its
validators check that the SHA-256 is exactly sixty-four hex
characters.

\subsection{Atomic JSON writing}
\begin{lstlisting}[language=python]
def write_json_model_atomic(
    path: Path,
    model: BaseModel,
    *,
    overwrite: bool = False,
) -> None:
\end{lstlisting}

The function writes a Pydantic model to disk atomically. It creates
a temporary file in the same directory, writes the model's JSON
representation with sorted keys and newline-terminated output,
\texttt{fsync}s the file, and then atomically renames the temporary
file over the target path. If any step fails, the temporary file is
removed before the exception propagates.

The \texttt{read\_json\_model()} function reads a JSON file and
validates it against a model type.

\subsection{Fingerprints}
The \texttt{artifacts/fingerprints.py} module provides three
streaming helpers:

\begin{itemize}[leftmargin=1.5em,itemsep=0.15em]
  \item \texttt{sha256\_file()} computes a file's SHA-256 by
    streaming it through a buffer, without loading it fully into
    memory;
  \item \texttt{file\_size\_bytes()} returns the exact size of an
    existing file;
  \item \texttt{count\_nonempty\_lines()} counts the non-empty lines
    in a file by streaming.
\end{itemize}

All three raise \texttt{FileNotFoundError} if the path does not
exist. They are used by the sidecar builders to construct manifests
without ever loading a large artifact into memory.

% =====================================================================
\section{Determinism and Reproducibility}\label{sec:determinism}

The architecture's determinism guarantees can be summarized at four
levels.

\subsection{Identity}
Every artifact has a content-addressable identity that is a
deterministic function of its inputs:

\begin{itemize}[leftmargin=1.5em,itemsep=0.15em]
  \item a parent's identity depends on the dataset, document ID,
    parser configuration hash, and parent index;
  \item a chunk's identity depends on the dataset, document ID,
    parent ID, chunker configuration hash, and parent-local chunk
    index;
  \item a Qdrant point's identity depends on the chunk ID;
  \item an embedding shard's identity depends on the dataset, model
    ID, embedding configuration hash, and shard index.
\end{itemize}

Because identities are deterministic, they are stable across runs,
across machines, and across re-executions with the same inputs. This
is what allows the pipeline to verify a resumed run's state by
comparing identity sets.

\subsection{Serialization}
JSON and JSONL artifacts are written with \texttt{sort\_keys=True}
and \texttt{separators=(",", ":")}. YAML artifacts are written with
\texttt{safe\_dump(sort\_keys=True)}. Pydantic models are serialized
in JSON mode with \texttt{exclude\_none=False}, so missing optional
fields appear as \texttt{null} rather than being omitted.

This combination means that two executions with the same inputs
produce byte-identical artifacts.

\subsection{Tie-breaking}
Wherever multiple items could be ranked equally by score, the
architecture breaks the tie by a deterministic secondary key:

\begin{itemize}[leftmargin=1.5em,itemsep=0.15em]
  \item parent aggregation: descending score, then best chunk
    rank, then parent ID;
  \item document resolution: descending score, then best parent
    rank, then document ID;
  \item RRF fusion (numpy path): descending RRF score, then
    ascending chunk index.
\end{itemize}

Tie-breaks are never left to chance. This is what makes the
architecture's rankings fully reproducible.

\subsection{Floating-point isolation}
The only source of nondeterminism in the pipeline is
floating-point arithmetic in the embedder. BGE-M3's FP16 forward
pass may produce slightly different vectors on different hardware
or in different library versions. This nondeterminism is isolated
to the embedding layer: everything downstream of the embedding
records is bit-exact given the same records.

% =====================================================================
\section{Package Structure and Entry Points}\label{sec:packages}

\subsection{Subpackage inventory}
Table~\ref{tab:packages} catalogs the subpackages of
\texttt{rag\_system} and their principal modules.

\begin{table}[h]
\centering
\caption{Subpackages of \texttt{rag\_system}.}
\label{tab:packages}
\small
\begin{tabular}{l l}
\toprule
\textbf{Subpackage} & \textbf{Modules} \\
\midrule
\texttt{artifacts} & fingerprints, identity, jsonl, manifest, manifest\_io, models, revision, sidecar \\
\texttt{canonical} & inspection, jsonl, loaders, models, paths, pilot\_adapter, pilot\_adapter\_persistence \\
\texttt{chunking} & config, hierarchical, huggingface, pipeline, tokenization, windows \\
\texttt{embedding} & catalog, config, identity, models, shards, synthetic \\
\texttt{evaluation} & clerk, identity, metrics, models, pilot\_persistence, pipeline, statistics \\
\texttt{generation} & (optional; disabled in retrieval-only execution) \\
\texttt{orchestration} & planner \\
\texttt{parsing} & blocks, config, deterministic, headings, pipeline, profiling, segmentation \\
\texttt{retrieval} & aggregation, document\_resolution, models, pilot\_persistence, qdrant \\
\texttt{storage} & models, qdrant \\
\bottomrule
\end{tabular}
\end{table}

\subsection{Entry points}
Table~\ref{tab:entrypoints} lists the entry points that external
code uses to invoke each stage of the pipeline.

\begin{table}[h]
\centering
\caption{Principal entry points.}
\label{tab:entrypoints}
\small
\begin{tabular}{l l}
\toprule
\textbf{Stage} & \textbf{Entry point} \\
\midrule
Canonical validation & \texttt{canonical.paths.resolve\_dataset\_paths} \\
Canonical streaming & \texttt{canonical.loaders.iter\_canonical\_documents/examples} \\
Parsing & \texttt{parsing.pipeline.parse\_canonical\_dataset} \\
Chunking & \texttt{chunking.pipeline.chunk\_parsed\_artifacts} \\
Embedding shard write & \texttt{embedding.shards.write\_embedding\_shard} \\
Qdrant collection & \texttt{storage.qdrant.create\_hybrid\_collection} \\
Qdrant ingestion & \texttt{storage.qdrant.ingest\_embedding\_records} \\
Dense retrieval & \texttt{retrieval.qdrant.retrieve\_dense} \\
Sparse retrieval & \texttt{retrieval.qdrant.retrieve\_sparse} \\
Hybrid RRF retrieval & \texttt{retrieval.qdrant.retrieve\_hybrid\_rrf} \\
Parent aggregation & \texttt{retrieval.aggregation.aggregate\_chunks\_to\_parents} \\
Document resolution & \texttt{retrieval.document\_resolution.resolve\_parents\_to\_documents} \\
Evaluation & \texttt{evaluation.pipeline.write\_retrieval\_evaluation} \\
\bottomrule
\end{tabular}
\end{table}

% =====================================================================
\section{Conclusion}

We have described the architecture and implementation of My\_RAG at
implementation granularity. The architecture is organized as a
sequence of nine subpackages, each with a well-defined
responsibility, a set of typed data models, and a set of invariants
that its layer is responsible for maintaining. The pipeline is
deterministic at every layer that produces artifacts, and the only
source of nondeterminism is isolated to the embedder. Every stage
publishes checksum-closed artifacts, and every stage is designed to
abort loudly when an invariant is violated rather than silently
propagate corrupted data.

The document is intended as a reference for readers who want to
understand the architecture from the source code without having to
read the source code directly. It is not a substitute for the
source, and it does not report the architecture's retrieval quality;
those results are documented separately.

% =====================================================================
\appendix

\section{Appendix A: Full Data Model Reference}

\subsection{A.1 --- Canonical models}
\begin{lstlisting}[language=python]
class CanonicalDocument(CanonicalBaseModel):
    doc_id: str
    dataset: str
    text_raw: str
    text_normalized: str
    title: str | None = None
    source_path: str | None = None
    source_row: int | str | None = None
    metadata: JsonObject = Field(default_factory=dict)

class CanonicalQuery(CanonicalBaseModel):
    query_id: str
    dataset: str
    question: str
    task_type: str | None = None
    split: str | None = None
    doc_ids: list[str] = Field(default_factory=list)
    metadata: JsonObject = Field(default_factory=dict)

class CanonicalGroundTruth(CanonicalBaseModel):
    ground_truth_type: GroundTruthType
    dataset: str
    query_id: str
    answers: list[Any] = Field(default_factory=list)
    relevant_ids: list[str] = Field(default_factory=list)
    choice: Any | None = None
    spans: list[Any] = Field(default_factory=list)
    metadata: JsonObject = Field(default_factory=dict)

class CanonicalExample(CanonicalBaseModel):
    query: CanonicalQuery
    ground_truth: CanonicalGroundTruth
    supplied_context: Any | None = None
    options: Any | None = None
    metadata: JsonObject = Field(default_factory=dict)
\end{lstlisting}

\subsection{A.2 --- Artifact models}
\begin{lstlisting}[language=python]
class OffsetSpan(ArtifactBaseModel):
    start: int = Field(ge=0)
    end: int = Field(ge=0)

class SemanticParent(ArtifactBaseModel):
    parent_id: str
    doc_id: str
    dataset: str
    parent_parent_id: str | None = None
    hierarchy_level: int = Field(ge=0)
    sibling_index: int = Field(ge=0)
    title: str | None = None
    text_normalized: str
    text_raw: str | None = None
    normalized_char_span: OffsetSpan
    raw_char_span: OffsetSpan | None = None
    token_span: OffsetSpan | None = None
    token_count: int | None = Field(default=None, ge=0)
    start_boundary: tuple[BoundaryEvidence, ...] = ()
    end_boundary: tuple[BoundaryEvidence, ...] = ()
    parser_id: str
    parser_version: str
    parser_config_hash: str
    metadata: JsonObject = Field(default_factory=dict)

class ParsedDocument(ArtifactBaseModel):
    artifact_type: str = "parsed_document"
    artifact_version: str = "1.0"
    doc_id: str
    dataset: str
    parser_id: str
    parser_version: str
    parser_config_hash: str
    canonical_text_normalized_length: int = Field(ge=0)
    canonical_text_raw_length: int = Field(ge=0)
    parents: tuple[SemanticParent, ...]
    source_fingerprint: str | None = None
    metadata: JsonObject = Field(default_factory=dict)

class ChildChunk(ArtifactBaseModel):
    artifact_type: str = "child_chunk"
    artifact_version: str = "1.0"
    chunk_id: str
    parent_id: str
    doc_id: str
    dataset: str
    chunk_index: int = Field(ge=0)
    parent_chunk_index: int = Field(ge=0)
    text_normalized: str
    text_raw: str | None = None
    normalized_char_span: OffsetSpan
    parent_normalized_char_span: OffsetSpan
    raw_char_span: OffsetSpan | None = None
    token_span: OffsetSpan | None = None
    token_count: int = Field(ge=0)
    parser_id: str
    parser_version: str
    parser_config_hash: str
    chunker_id: str
    chunker_version: str
    chunker_config_hash: str
    tokenizer_id: str
    chunk_size: int = Field(gt=0)
    chunk_overlap: int = Field(ge=0)
    metadata: JsonObject = Field(default_factory=dict)
\end{lstlisting}

\subsection{A.3 --- Retrieval models}
\begin{lstlisting}[language=python]
class RetrievedChunk(ArtifactBaseModel):
    chunk_id: str
    parent_id: str
    doc_id: str
    dataset: str
    rank: int = Field(ge=1)
    score: float
    text_normalized: str
    retrieval_method: str
    payload: dict[str, object] = Field(default_factory=dict)

class RetrievedParent(ArtifactBaseModel):
    parent_id: str
    doc_id: str
    dataset: str
    rank: int = Field(ge=1)
    score: float
    best_chunk_rank: int = Field(ge=1)
    matched_chunk_count: int = Field(ge=1)
    chunk_ids: tuple[str, ...]
    child_scores: tuple[float, ...]

class RetrievedDocument(ArtifactBaseModel):
    query_id: str
    doc_id: str
    dataset: str
    rank: int = Field(ge=1)
    score: float
    best_parent_rank: int = Field(ge=1)
    contributing_parent_ids: tuple[str, ...]
    contributing_chunk_ids: tuple[str, ...]
    metadata: dict[str, object] = Field(default_factory=dict)
\end{lstlisting}

\subsection{A.4 --- Embedding models}
\begin{lstlisting}[language=python]
class SparseVector(ArtifactBaseModel):
    indices: tuple[int, ...]
    values: tuple[float, ...]

class ChunkEmbeddingRecord(ArtifactBaseModel):
    artifact_type: str = "chunk_embedding"
    artifact_version: str = "1.0"
    chunk_id: str
    parent_id: str
    doc_id: str
    dataset: str
    model_id: str
    model_revision: str
    tokenizer_id: str
    embedding_config_hash: str
    dense_vector: tuple[float, ...] | None = None
    sparse_vector: SparseVector | None = None
    dense_dimension: int | None = Field(default=None, gt=0)
    dense_precision: VectorPrecision | None = None
    dense_normalized: bool | None = None
    source_chunk_sha256: str | None = None

class EmbeddingShardHeader(ArtifactBaseModel):
    artifact_type: str = "embedding_shard_header"
    artifact_version: str = "1.0"
    shard_id: str
    shard_index: int = Field(ge=0)
    dataset: str
    model_id: str
    model_revision: str
    tokenizer_id: str
    embedding_config_hash: str
    embedding_kind: EmbeddingKind
    dense_dimension: int | None = Field(default=None, gt=0)
    dense_precision: VectorPrecision | None = None
    dense_normalized: bool | None = None
    maximum_records: int = Field(gt=0)
\end{lstlisting}

\section{Appendix B: Configuration Reference}

\subsection{B.1 --- DeterministicParserConfig}
\begin{lstlisting}[language=python]
class DeterministicParserConfig(BaseModel):
    model_config = ConfigDict(extra="forbid", frozen=True)

    parser_id: str = "deterministic_structural_parser"
    parser_version: str = "0.3.0"

    detect_numbered_headings: bool = True
    detect_named_headings: bool = True
    detect_uppercase_headings: bool = True

    maximum_heading_characters: int = Field(default=120, ge=10)
    maximum_heading_words: int = Field(default=14, ge=1)
    minimum_heading_alphabetic_characters: int = Field(
        default=4, ge=1,
    )
    maximum_heading_terminal_punctuation: int = Field(
        default=1, ge=0,
    )

    reject_equation_like_blocks: bool = True
    reject_url_like_blocks: bool = True
    reject_numeric_table_fragments: bool = True

    minimum_parent_characters: int = Field(default=20, ge=1)
    merge_small_structural_groups: bool = True
    preserve_raw_text_when_identical: bool = True
\end{lstlisting}

\subsection{B.2 --- ChildChunkerConfig}
\begin{lstlisting}[language=python]
class ChildChunkerConfig(BaseModel):
    model_config = ConfigDict(extra="forbid", frozen=True)

    chunker_id: str = "hierarchical_token_chunker"
    chunker_version: str = "0.3.0"
    chunk_size: int = Field(default=256, gt=0)
    chunk_overlap: int = Field(default=32, ge=0)
    minimum_final_chunk_tokens: int | None = Field(
        default=None, gt=0,
    )
    preserve_parent_boundaries: bool = True
    include_empty_parents: bool = False
\end{lstlisting}

\subsection{B.3 --- EmbeddingConfig}
\begin{lstlisting}[language=python]
class EmbeddingConfig(BaseModel):
    model_config = ConfigDict(extra="forbid", frozen=True)

    embedder_id: str = "bge_m3_embedder"
    embedder_version: str = "0.1.0"
    model_id: str = "BAAI/bge-m3"
    model_revision: str
    tokenizer_id: str
    embedding_kind: EmbeddingKind = EmbeddingKind.HYBRID
    dense_dimension: int | None = Field(default=1024, gt=0)
    dense_precision: VectorPrecision | None = VectorPrecision.FLOAT32
    dense_normalized: bool | None = True
    device: str = "cpu"
    use_fp16: bool = False
    batch_size: int = Field(default=2, gt=0)
    maximum_model_input_tokens: int = Field(default=8192, gt=0)
    shard_maximum_records: int = Field(default=10000, gt=0)
\end{lstlisting}

\section{Appendix C: Determinism Guarantees Summary}

\begin{center}
\small
\begin{tabular}{p{4cm} p{10.5cm}}
\toprule
\textbf{Guarantee} & \textbf{Implementation} \\
\midrule
Content-addressable IDs & \texttt{build\_parent\_id}, \texttt{build\_chunk\_id}, \texttt{build\_embedding\_shard\_id}, \texttt{build\_qdrant\_point\_id} \\
Sorted-key JSON & \texttt{json.dumps(..., sort\_keys=True, separators=(",", ":"))} \\
Atomic writes & \texttt{write\_json\_model\_atomic}, \texttt{AtomicJsonlWriter} \\
Deterministic tie-breaks & Explicit secondary keys in aggregation and resolution \\
Parser determinism & Fixed regex rules, no randomness \\
Chunker determinism & Fixed token-window algorithm, no randomness \\
Fingerprint chain & Every manifest records the SHA-256 of its inputs \\
Floating-point isolation & Nondeterminism confined to embedder FP16 arithmetic \\
\bottomrule
\end{tabular}
\end{center}

\end{document}